%% kapa: exam=1 round=22 year=2011 period=1 form=A subject=civil_law subject_ko=민법 seq=37 number=37
%% source: Q-Net 감정평가사 기출문제 (cst00309) artlSeq=5036790 file=extracted/1008/5036790/2011년 제22회 감정평가사 1차시험 기출문제/1교시.hwp
%% conversion: kordoc 4.24.51 -> tools/gen1.py (auto); review=proofread-sonnet; flags=none
\begin{question}{37}
지역권에 관한 설명으로 옳은 것은?
\begin{choices}
  \item 지역권은 시효취득 할 수 없다.
  \item 지역권은 요역지와 분리하여 양도할 수 있다.
  \item 지역권은 다른 약정이 없는 한 요역지소유권에 부종하여 이전한다.
  \item 공유자의 1인이 지역권을 취득한 때에도 다른 공유자는 이를 취득하지 못한다.
  \item 요역지가 수인의 공유인 경우에 그 1인에 의한 지역권 소멸시효의 중단은 다른 공유자를 위하여 효력이 없다.
\end{choices}
\end{question}
\answer{3}
